\documentclass[12pt,a4paper]{book}
\usepackage{fontspec}
\setmainfont{Times New Roman}
\usepackage[margin=1in]{geometry}
\usepackage[hidelinks]{hyperref}
\usepackage{pdfpages}
\usepackage{setspace}

\begin{document}
\frontmatter
\begin{titlepage}
\thispagestyle{empty}
\begin{center}
{\Large Ontology of Continua}\\[0.6cm]
{\huge \textbf{Core 1.3}}\\[0.3cm]
{\Large Русское проводящее издание монографии}\\[1.2cm]

{\Large Alexander Yashin}\\[0.2cm]
{\normalsize Independent Researcher}\\[0.2cm]
{\normalsize ORCID 0009-0008-6166-0914}\\[1.0cm]

{\normalsize Version 1.3}\\[0.2cm]
{\normalsize April 13, 2026}\\[1.2cm]

\begin{minipage}{0.88\textwidth}
\small
\textbf{Статус издания.} Это русское проводящее издание основной монографии OC Core 1.3.
Оно даёт локализованную точку входа, объясняет архитектуру доказательного корпуса,
очерчивает границы утверждений и затем включает полный английский master-volume как
каноническое theorem-bearing тело.

\medskip

\textbf{Почему так.} На текущем релизном шаге канонический доказательный корпус,
собранный непосредственно из публичного LaTeX-источника, остаётся английским master.
Русское издание поэтому делает главное: помогает русскоязычному читателю правильно
войти в работу, понять её структуру, статус и смысл доказательств, а затем читать сам
полный корпус без потери научной точности.
\end{minipage}

\vfill

{\small Open foundation route: github.com/alexanderyashin/ontology-of-continua-core-main}\\
{\small Archival predecessor witness: Zenodo DOI 10.5281/zenodo.17903912}
\end{center}
\end{titlepage}

\chapter*{Аннотация}
\addcontentsline{toc}{chapter}{Аннотация}
\onehalfspacing
Core 1.3 представляет собой source-first реконструкцию и публикационное ужесточение
монографического корпуса Ontology of Continua. Научная задача этого релиза двойная:
сохранить полный масштаб formal shell, inherited from Core 1.2, и одновременно сделать
явными source witness discipline, theorem-fate separation, chronology-aware repair notes,
границы утверждений и outward-ready scholarly packaging.

Это издание не заменяет английскую theorem-bearing монографию и не маскирует её под
сокращённый packet. Его функция состоит в том, чтобы русскоязычный читатель получил
ясный narrative entry point: что именно здесь доказывается, как устроена архитектура
теорем и лемм, почему зафиксирована иерархия `K0…K12`, где заканчивается положительное
тело доказанной теории и где начинаются bounded interpretations, cross-domain projections
и дальнейшие программы проверки.

\tableofcontents

\mainmatter

\chapter{Как читать это издание}
Монография OC Core 1.3 рассчитана одновременно на три читательских режима.

\begin{enumerate}
\item \textbf{Редактор и рецензент} должны быстро увидеть, что является theorem-bearing
телом, где лежат proofs, какие ограничения задекларированы явно и какие утверждения
вынесены за пределы положительного корпуса.
\item \textbf{Формальный читатель} должен иметь непрерывный доступ к аксиомам,
определениям, теоремам, леммам, формулам и приложениями без необходимости восстанавливать
структуру из внутренних ledger-файлов.
\item \textbf{Неспециалист, но серьёзный научный читатель} должен получить понятную
точку входа, narrative spine и пояснения, зачем каждая большая глава вообще существует.
\end{enumerate}

Поэтому это проводящее издание делает две вещи. Сначала оно даёт локализованную карту
чтения. Затем оно включает полный английский master-volume, который и является
каноническим доказательным телом релиза.

\chapter{Что именно доказывается в Core 1.3}
В положительном теле релиза центральной остаётся bounded grammar collapse, residue and
rebirth. Это означает следующее.

\begin{itemize}
\item При потере непустой допустимой реализации в текущем режиме вложения и порога
система переходит в lawful collapse state.
\item Постколлапсный след может сохраняться как residue.
\item Более позднее допустимое продолжение не считается сохранением той же самой
идентичности автоматически; оно классифицируется как rebirth и требует собственной
идентификационной дисциплины.
\end{itemize}

Core 1.3 усиливает не столько риторику, сколько дисциплину. Он отделяет surviving
source-native kernel от исключённых theorem families, явным образом фиксирует theorem fate,
добавляет source-audit layers и не позволяет summary-слоям выдать внутренний оптимизм за
публичное доказательство.

\chapter{Почему иерархия фиксируется как \texorpdfstring{$K0 \dots K12$}{K0...K12}}
Core 1.3 устраняет прежнюю двусмысленность, при которой текстами местами говорилось
`K0–K10`, тогда как реальные source surfaces и dependency-bearing modules уже работали с
более длинной train-структурой. В финальном релизном doctrine hierarchy is fixed as
`K0…K12`.

\begin{itemize}
\item `K0…K10` продолжают покрывать базовую structural ladder и механизм роста
выразительности.
\item `K11` нужен для рекурсивных meta-theoretical axes, где описание уже управляет не
только объектами и переходами, но и перестройкой самих explanatory regimes.
\item `K12` нужен для global coherence axes, где возникает вопрос не о локальной
выразительности отдельного уровня, а о согласовании всей train-структуры как единого
объекта.
\end{itemize}

Именно поэтому добавленные уровни не являются декоративным счётом. Они нужны, чтобы
иерархия не обрывалась в точке, где теория уже фактически обсуждает meta-level closure и
whole-structure coherence.

\chapter{Карта полной монографии}
Полный английский master-volume, включённый далее, следует читать как opus-magnus source.
Его структура в крупном масштабе такова.

\begin{enumerate}
\item Front matter задаёт publication role, reader promise, abstract и общую цель работы.
\item Основные главы вводят формальную модель, embedding spaces, axioms, operator logic и
теоремный корпус.
\item Отдельные главы разворачивают collapse / death / residue / rebirth semantics,
cross-level constraints и domain-facing interpretive boundary.
\item Reader-guide, source-audit и foundational-consistency chapters добавляют именно то,
чего не хватало для Core 1.3: понятную didactic spine и явную publication discipline.
\item Appendix layers удерживают provenance, source audit, extraction boundary between
master monograph and journal core, а также дополнительные support materials.
\end{enumerate}

\chapter{Граница claims и что именно не утверждается}
Scientific value этой монографии держится не на риторическом максимализме, а на строгом
разделении утверждений.

\begin{itemize}
\item \textbf{Axiomatic statements} задают базовые формальные посылки.
\item \textbf{Derived statements} входят в theorem-bearing корпус только там, где
dependency chain и proof route прослеживаются явно.
\item \textbf{Structural interpretations} допустимы как объясняющие слои, но не
автоматически как новые теоремы.
\item \textbf{Cross-domain projections} служат картой дальнейшей работы, а не публичным
доказательством замыкания всех областей сразу.
\end{itemize}

Поэтому в Core 1.3 сознательно не делаются заявления, которые шире подтверждённого
proof surface. Ambition framework сохраняется, но public wording остаётся sober,
exact and conventionally scientific.

\chapter{Как соотносится journal core}
Journal core не заменяет монографию. Он является traceable extraction from the master.
Его задача --- дать редактору и рецензенту shorter editorial route, не теряя канонической
связи с большим доказательным корпусом.

Правильный порядок чтения таков.

\begin{enumerate}
\item Если нужен полный доказательный и дидактический контекст, начинать надо с master
monograph.
\item Если нужен компактный editorial route для журнала, можно начинать с journal core,
но всё равно рассматривать его как extract, а не как отдельную теорию.
\item Все companion packets должны толковаться как support layers вокруг этих двух
science artifacts, а не как замена самой рукописи.
\end{enumerate}

\chapter{Канонический английский корпус}
Далее следует полный английский master-volume, built directly from the public LaTeX source
corpus. Именно он является canonical theorem-bearing body релиза Core 1.3.

\cleardoublepage
\includepdf[pages=-]{build_oc_core_1_3_master/oc_core_1_3_master_monograph.pdf}

\end{document}
